\begin{figure}[t]
\centering
\begingroup
\definecolor{overviewblue}{HTML}{255C8A}
\definecolor{overviewgreen}{HTML}{237A65}
\begin{tikzpicture}[
  font=\small,
  source/.style={draw,rounded corners,minimum width=17mm,
                 minimum height=8mm,fill=black!3},
  chosen/.style={source,draw=overviewblue,fill=overviewblue!10,thick},
  value/.style={draw,circle,minimum size=7mm,inner sep=0pt,fill=white},
  affected/.style={value,draw=overviewblue,fill=overviewblue!10},
  link/.style={-{Stealth[length=1.6mm]},draw=black!45},
  chosenlink/.style={link,draw=overviewblue,thick},
  note/.style={font=\footnotesize,align=center}]

\node[draw=black!20,rounded corners,fill=black!3,minimum width=8cm,
      minimum height=13mm,align=center] at (4,6.9)
  {\textbf{Question: which source tuples must we delete?}\\[2pt]
   Fewest deletions, within budget $k$, causing change $\ge\theta$};

\node[font=\footnotesize\bfseries,align=center] at (4,5.85)
  {Running example: $Q=\agg{SUM}$, initial answer $15$};
\node[note] at (1.55,5.3) {Deletable source tuples};
\node[note] at (6.25,5.3) {Join contributions};

\node[chosen] (r1) at (1.55,4.4) {$r_1$};
\node[source] (r2) at (1.55,3.15) {$r_2$};
\node[affected] (v4) at (6.25,4.75) {$4$};
\node[affected] (v6) at (6.25,3.95) {$6$};
\node[value] (v5) at (6.25,3.15) {$5$};
\draw[chosenlink] (r1.east)--(v4.west);
\draw[chosenlink] (r1.east)--(v6.west);
\draw[link] (r2.east)--(v5.west);
\node[note,text=overviewblue] at (3.85,2.55)
  {Deleting $r_1$ removes two contributions: $4+6$};

\node[draw=black!20,rounded corners,minimum width=8cm,
      minimum height=12mm,align=center] at (4,1.55)
  {$k=1$: at most one source deletion\\[2pt]
   $\theta=8$: lower the sum from $15$ to at most $7$};

\node[draw=overviewgreen,rounded corners,fill=overviewgreen!7,
      minimum width=38mm,minimum height=14mm,align=center] (accepted) at (1.95,0.1)
  {\textbf{Delete $r_1$: succeeds}\\[2pt]
   $15\to5$; loss $10\ge8$};
\node[draw=black!30,rounded corners,fill=black!2,
      minimum width=38mm,minimum height=14mm,align=center] at (6.05,0.1)
  {\textbf{Delete $r_2$: fails}\\[2pt]
   $15\to10$; loss $5<8$};

\node[draw=overviewblue,rounded corners,fill=overviewblue!8,
      minimum width=8cm,minimum height=14mm,align=center] (answer) at (4,-1.6)
  {\textbf{Minimum witness: $S^\star=\{r_1\}$, $|S^\star|=1\le k$}\\[2pt]
   Zero deletions cannot reach the threshold.};
\draw[link,draw=overviewgreen] (accepted.south) -- (accepted.south |- answer.north);
\end{tikzpicture}
\endgroup
\caption{The paper's objective, illustrated for SUM. Minimize the number of
permitted source deletions subject to $|S|\le k$ and $\Delta(S)\ge\theta$.
The budget $k$ counts source tuples; the threshold $\theta$ measures the
answer change. If no valid set fits the budget, report that outcome.}
\label{fig:overview}
\Description{An overview beside the abstract explains minimum k-instability
witnesses with a SUM example. Source tuple r1 contributes join values 4 and
6, while r2 contributes 5, giving an initial sum of 15. With at most one
source deletion and a required drop of 8, deleting r1 leaves 5 and succeeds;
deleting r2 leaves 10 and fails. The minimum witness is the single tuple r1.
This separates the deletion budget k from the answer-change threshold theta.}
\end{figure}
